% =============================================================================
%                       CHAPTER 4: IMPLEMENTATION
% =============================================================================
\chapter{Implementation}

\section{Tools and Techniques}

\subsection{Hardware Environment}

\begin{table}[H]
\centering
\caption{Hardware configuration for experiments.}
\label{tab:hardware_config}
\begin{tabular}{ll}
\toprule
\textbf{Component} & \textbf{Specification} \\
\midrule
GPU & NVIDIA Tesla T4 (15.6 GB VRAM) \\
CPU & Intel Xeon (2 cores) \\
RAM & 13 GB System Memory \\
Platform & Kaggle Notebooks \\
CUDA Version & 12.x \\
Persistent Storage & \texttt{/kaggle/working/} (saved as output) \\
\bottomrule
\end{tabular}
\end{table}

The NVIDIA Tesla T4 GPU, based on the Turing architecture, provides 15.6 GB of VRAM with support for mixed-precision (FP16) and INT8/INT4 operations. While not sufficient to load Qwen2.5-VL at full precision (which requires approximately 14 GB in float16), the T4 supports 4-bit quantization through the BitsAndBytes library, enabling deployment of the 7B-parameter model within the available memory budget.

\subsection{Software Stack}

\begin{table}[H]
\centering
\caption{Software libraries and their roles in the experimental pipeline.}
\label{tab:software_stack}
\begin{tabularx}{\textwidth}{l l X}
\toprule
\textbf{Library} & \textbf{Version} & \textbf{Purpose} \\
\midrule
Python & 3.12.13 & Programming language \\
PyTorch & 2.x (Latest) & Deep learning framework, CUDA backend, custom NN classifiers \\
Transformers & Latest & Hugging Face model hub, V-JEPA 2 and Qwen2.5-VL loading \\
BitsAndBytes & Latest & 4-bit NF4 quantization for Qwen2.5-VL \\
scikit-learn & Latest & Classical ML classifiers, PCA, metrics \\
OpenCV & Latest & Video frame extraction (\texttt{cv2.VideoCapture}) \\
yt-dlp & Latest & YouTube video downloading with temporal trimming \\
huggingface\_hub & Latest & HMDB51 dataset download (\texttt{hf\_hub\_download}) \\
NumPy & Latest & Numerical operations, frame indexing \\
Pillow (PIL) & Latest & Image format conversion \\
\bottomrule
\end{tabularx}
\end{table}

\subsection{Model Specifications}

\begin{table}[H]
\centering
\caption{Architectural specifications of V-JEPA 2 and Qwen2.5-VL.}
\label{tab:model_specs}
\begin{tabularx}{\textwidth}{l X X}
\toprule
\textbf{Specification} & \textbf{V-JEPA 2} & \textbf{Qwen2.5-VL} \\
\midrule
Parameters & 654M & 7B \\
Architecture & ViT-Huge & Qwen2.5 + ViT \\
Paradigm & Self-supervised (JEPA) & Autoregressive VLM \\
Pretraining & Masked latent prediction & Image-text pairs + instruction tuning \\
Input Frames & 8 & 4 \\
Input Resolution & $224 \times 224$ & Dynamic resolution \\
Feature Dimension & 1,280 & 3,584 \\
Precision & float16 + SDPA & 4-bit NF4 \\
Modality & Video only & Vision + Language \\
Hugging Face ID & \nolinkurl{facebook/vjepa2-vith-fpc64-256} & \nolinkurl{Qwen/Qwen2.5-VL-7B-Instruct} \\
\bottomrule
\end{tabularx}
\end{table}

\section{Implementation Details}

\subsection{FineGym Data Pipeline}

The FineGym pipeline downloads clips from YouTube using \texttt{yt-dlp} with temporal trimming. The annotation JSON provides YouTube video IDs with event timestamps, which are used to extract precise event segments. The following code snippet illustrates the download function:

\begin{lstlisting}[caption={FineGym clip download with \texttt{yt-dlp} temporal trimming.}, label={lst:fg_download}]
def download_clip(yt_id, start, end, out_path, max_retries=2):
    """Download a clip from YouTube using yt-dlp."""
    if os.path.exists(out_path) and os.path.getsize(out_path) > 1000:
        return True # Already downloaded
    url = f'https://www.youtube.com/watch?v={yt_id}'
    for attempt in range(max_retries):
        try:
            cmd = [
                'yt-dlp', '--quiet', '--no-warnings',
                '--download-sections', f'*{start}-{end}',
                '-f', 'worst[ext=mp4]/worst',
                '--force-keyframes-at-cuts',
                '-o', out_path, url
            ]
            subprocess.run(cmd, capture_output=True, text=True, timeout=120)
            if os.path.exists(out_path) and os.path.getsize(out_path) > 500:
                return True
        except (subprocess.TimeoutExpired, Exception):
            pass
    return False
\end{lstlisting}

Key implementation decisions for the FineGym pipeline:
\begin{itemize}[leftmargin=*]
    \item \textbf{Lowest Quality Selection:} The \texttt{worst[ext=mp4]/worst} format selector downloads the lowest available quality, minimizing download time and storage while preserving sufficient visual information for frame extraction at $224 \times 224$.
    \item \textbf{Temporal Trimming:} The \nolinkurl{--download-sections} flag with \nolinkurl{--force-keyframes-at-cuts} ensures precise temporal extraction matching the annotation timestamps.
    \item \textbf{Resume Support:} A JSON log file tracks download status for each clip, enabling the pipeline to resume after interruption without re-downloading successful clips.
    \item \textbf{Balanced Sampling:} 300 clips per class are sampled to ensure balanced representation, yielding 1,200 clips across the four event classes.
\end{itemize}

\subsection{HMDB51 Data Pipeline}

HMDB51 is downloaded from Hugging Face (\texttt{jili5044/hmdb51}) as a 2.12 GB zip archive using \texttt{hf\_hub\_download}. The implementation bypasses the standard Hugging Face \texttt{datasets} library, which encountered schema validation errors, by directly downloading and extracting the raw zip archive:

\begin{lstlisting}[caption={HMDB51 download and local indexing pipeline.}, label={lst:hmdb_download}]
from huggingface_hub import hf_hub_download

# Download raw zip directly (bypasses schema errors)
archive_path = hf_hub_download(
    repo_id="jili5044/hmdb51",
    filename="hmdb51.zip",
    repo_type="dataset",
    local_dir=download_dir
)

# Extract and index videos by class folder
with zipfile.ZipFile(archive_path, 'r') as zip_ref:
    zip_ref.extractall(extract_dir)

# Build sample list from directory structure
ds = []
for root, dirs, files in os.walk(extract_dir):
    for file in files:
        if file.endswith(('.mp4', '.avi', '.mkv')):
            video_path = os.path.join(root, file)
            label_name = os.path.basename(root)
            ds.append({'video_path': video_path, 'label': label_name})
\end{lstlisting}

After extraction, video files are indexed by walking the directory tree, with action class labels inferred from parent folder names. All 6,754 videos (51 classes) are successfully loaded with zero errors.

\subsection{Feature Extraction Implementation}

Both V-JEPA 2 and Qwen2.5-VL are loaded sequentially (one at a time) to manage GPU memory constraints within the 15.6 GB VRAM budget.

\subsubsection{V-JEPA 2 Extraction}

\begin{lstlisting}[caption={V-JEPA 2 feature extraction with checkpointing.}, label={lst:vjepa_extraction}]
from transformers import VJEPA2Model, VJEPA2VideoProcessor

vproc = VJEPA2VideoProcessor.from_pretrained(
    'facebook/vjepa2-vith-fpc64-256')
vmodel = VJEPA2Model.from_pretrained(
    'facebook/vjepa2-vith-fpc64-256',
    torch_dtype=torch.float16,
    attn_implementation='sdpa',
).to(device).eval()

# Feature extraction with checkpointing
for i in range(start_idx, total):
    frames = load_frames(clip['path'], 8, 224)
    inputs = vproc(frames, return_tensors='pt')
    inputs = {k: v.to(device) for k, v in inputs.items()}
    with torch.no_grad():
        with torch.amp.autocast('cuda'):
            out = vmodel(**inputs)
            feat = out.last_hidden_state.mean(dim=1).cpu().float().squeeze(0)
    features.append(feat)
\end{lstlisting}

\subsubsection{Qwen2.5-VL Extraction}

\begin{lstlisting}[caption={Qwen2.5-VL 4-bit quantized feature extraction.}, label={lst:qwen_extraction}]
from transformers import (
    Qwen2_5_VLForConditionalGeneration,
    AutoProcessor, BitsAndBytesConfig)

bnb = BitsAndBytesConfig(
    load_in_4bit=True,
    bnb_4bit_compute_dtype=torch.float16,
    bnb_4bit_quant_type='nf4',
    bnb_4bit_use_double_quant=True)

qmodel = Qwen2_5_VLForConditionalGeneration.from_pretrained(
    'Qwen/Qwen2.5-VL-7B-Instruct',
    quantization_config=bnb,
    device_map='auto',
    torch_dtype=torch.float16)
qproc = AutoProcessor.from_pretrained(
    'Qwen/Qwen2.5-VL-7B-Instruct')

# Chat template construction for feature extraction
messages = [{'role': 'user', 'content': [
    {'type': 'video', 'video': frames},
    {'type': 'text', 'text': 'Describe.'},
]}]
text = qproc.apply_chat_template(
    messages, tokenize=False, add_generation_prompt=True)
inputs = qproc(text=[text], videos=[frames],
               return_tensors='pt', padding=True)

with torch.no_grad():
    out = qmodel(**inputs, output_hidden_states=True)
    feat = out.hidden_states[-1].mean(dim=1).cpu().float().squeeze(0)
\end{lstlisting}

\subsection{Neural Network Classifier Implementation}

In addition to the three sklearn classifiers, five neural network architectures are implemented for frozen-feature classification. All share a common training harness with advanced regularization techniques. The following code snippet illustrates the ResMLP architecture and the training loop with label smoothing and mixup:

\begin{lstlisting}[caption={ResMLP architecture with residual blocks.}, label={lst:resmlp_arch}]
class ResBlock(nn.Module):
    def __init__(self, dim, dropout=0.2):
        super().__init__()
        self.net = nn.Sequential(
            nn.Linear(dim, dim),
            nn.BatchNorm1d(dim), nn.GELU(),
            nn.Dropout(dropout),
            nn.Linear(dim, dim),
            nn.BatchNorm1d(dim),
        )
        self.act = nn.GELU()

    def forward(self, x):
        return self.act(x + self.net(x))

class ResMLP(nn.Module):
    def __init__(self, in_dim, num_classes, hidden=256, n_blocks=3):
        super().__init__()
        self.proj = nn.Sequential(
            nn.Linear(in_dim, hidden),
            nn.BatchNorm1d(hidden), nn.GELU())
        self.blocks = nn.Sequential(
            *[ResBlock(hidden) for _ in range(n_blocks)])
        self.head = nn.Linear(hidden, num_classes)

    def forward(self, x):
        return self.head(self.blocks(self.proj(x)))
\end{lstlisting}

\begin{lstlisting}[caption={Training loop with label smoothing, mixup augmentation, and cosine LR scheduling.}, label={lst:training_loop}]
LABEL_SMOOTHING = 0.1

def mixup_data(x, y, alpha=0.2):
    """Feature-space mixup augmentation."""
    lam = np.random.beta(alpha, alpha)
    lam = max(lam, 1 - lam)
    idx = torch.randperm(x.size(0))
    return lam * x + (1-lam) * x[idx], y, y[idx], lam

def train_model(model, X_tr, y_tr, X_ev, y_ev, epochs=300, lr=1e-3, patience=30):
    opt = torch.optim.AdamW(model.parameters(), lr=lr, weight_decay=1e-4)
    scheduler = torch.optim.lr_scheduler.CosineAnnealingLR(opt, T_max=epochs)
    best_acc, best_state, wait = 0, None, 0

    for epoch in range(epochs):
        model.train()
        for xb, yb in dataloader:
            xb_mix, ya, yb_mix, lam = mixup_data(xb, yb)
            logits = model(xb_mix)
            loss = lam * F.cross_entropy(logits, ya, label_smoothing=LABEL_SMOOTHING) + \
                   (1 - lam) * F.cross_entropy(logits, yb_mix, label_smoothing=LABEL_SMOOTHING)
            opt.zero_grad()
            loss.backward()
            opt.step()
        scheduler.step()

        # Early stopping on eval accuracy
        model.eval()
        with torch.no_grad():
            acc = (model(X_ev).argmax(1) == y_ev).float().mean().item()
        if acc > best_acc:
            best_acc = acc
            best_state = model.state_dict()
            wait = 0
        else:
            wait += 1
            if wait >= patience: break
    model.load_state_dict(best_state)
    return model
\end{lstlisting}

\subsection{HAF-STD Net Classifier Implementation}

While baseline classifiers operate on temporally pooled, PCA-reduced single-layer representations, the proposed \textbf{HAF-STD Net} (Hierarchical Adaptive Fusion with Spectral-Temporal Decomposition) directly processes 4D multi-layer spatiotemporal tensors $(B, K=4, T, D=1280)$ extracted from intermediate layers of V-JEPA~2. Listing~\ref{lst:hafstd_arch_impl} provides the complete PyTorch implementation of the Haar discrete wavelet transform module and the HAF-STD Net architecture configured for the 51-class HMDB51 action recognition benchmark (from \texttt{newhmdb.ipynb}).

\begin{lstlisting}[caption={HAF-STD Net architecture and spectral-temporal decomposition (HMDB51 configuration).}, label={lst:hafstd_arch_impl}]
class HaarDWT1D(nn.Module):
    """1D Haar Discrete Wavelet Transform along temporal axis.
    Decomposes temporal signal into low-frequency (appearance/trend)
    and high-frequency (motion/change) components.
    """
    def forward(self, x):
        """x: (B, T, D) -> (low: (B, T//2, D), high: (B, T//2, D))"""
        if x.shape[1] % 2 != 0:
            x = F.pad(x, (0, 0, 0, 1))  # pad temporal dim if odd
        x_even = x[:, 0::2, :]  # Even-indexed temporal positions
        x_odd = x[:, 1::2, :]   # Odd-indexed temporal positions
        low = (x_even + x_odd) * 0.5    # Approximation = appearance/trend
        high = (x_even - x_odd) * 0.5   # Detail = motion/change dynamics
        return low, high


class HAFSTD_Net(nn.Module):
    """Hierarchical Adaptive Fusion with Spectral-Temporal Decomposition.
    Novel architecture for V-JEPA 2.1 action classification on HMDB51 (51 classes).
    """
    def __init__(self, embed_dim=1280, num_layers=4, num_classes=51,
                 hidden=256, dropout=0.3):
        super().__init__()
        self.num_layers = num_layers
        # 11 streams: 4 appearance, 4 motion, 3 depth velocity
        self.num_streams = (num_layers * 2) + (num_layers - 1)  # 4*2 + 3 = 11
        self.dwt = HaarDWT1D()

        # Dedicated projection heads (11 total)
        self.stream_projections = nn.ModuleList([
            nn.Sequential(
                nn.Linear(embed_dim, hidden),
                nn.LayerNorm(hidden),
                nn.GELU(),
                nn.Dropout(dropout),
            )
            for _ in range(self.num_streams)
        ])

        # Adaptive routing network over 11 depth x frequency streams
        self.router = nn.Sequential(
            nn.Linear(self.num_streams * hidden, hidden),
            nn.GELU(),
            nn.Dropout(dropout),
            nn.Linear(hidden, self.num_streams),
        )

        # Classification head with cosine similarity
        self.final_proj = nn.Sequential(
            nn.Linear(hidden, hidden),
            nn.LayerNorm(hidden),
            nn.GELU(),
            nn.Dropout(dropout),
        )
        self.class_weights = nn.Parameter(torch.randn(num_classes, hidden))
        self.scale = nn.Parameter(torch.tensor(16.0))

    def forward(self, x):
        """x: (B, num_layers=4, T, D=1280)"""
        B = x.shape[0]
        streams = []

        # 1. Appearance & Motion streams (4 layers x 2 = 8 streams)
        for k in range(self.num_layers):
            layer_feat = x[:, k, :, :]  # (B, T, D)
            low, high = self.dwt(layer_feat)
            appearance = low.mean(dim=1)  # low-freq = appearance
            motion = high.mean(dim=1)      # high-freq = motion dynamics
            streams.append(self.stream_projections[k * 2](appearance))
            streams.append(self.stream_projections[k * 2 + 1](motion))

        # 2. Depth Representation Velocity streams (3 inter-layer deltas)
        for k in range(self.num_layers - 1):
            next_layer = x[:, k + 1, :, :].mean(dim=1)
            curr_layer = x[:, k, :, :].mean(dim=1)
            velocity = next_layer - curr_layer
            streams.append(self.stream_projections[8 + k](velocity))

        # Stack all 11 streams: (B, 11, hidden)
        all_streams = torch.stack(streams, dim=1)

        # 3. Adaptive Routing
        flat = all_streams.reshape(B, -1)
        routing_weights = torch.softmax(self.router(flat), dim=-1)
        fused = (all_streams * routing_weights.unsqueeze(-1)).sum(dim=1)

        # 4. Cosine Similarity Classification
        fused = self.final_proj(fused)
        fused_norm = F.normalize(fused, p=2, dim=-1)
        weight_norm = F.normalize(self.class_weights, p=2, dim=-1)
        logits = self.scale * fused_norm @ weight_norm.T
        return logits
\end{lstlisting}

To train HAF-STD Net on the 4D feature tensors without collapsing intermediate layer information, the training pipeline generalizes feature-space mixup to multi-layer temporal tensors:

\begin{lstlisting}[caption={4D Tensor Mixup and Training Harness for HAF-STD Net on HMDB51.}, label={lst:hafstd_train_impl}]
def train_hafstd_model(model, X_tr_4d, y_tr, X_ev_4d, y_ev,
                       epochs=300, lr=1e-3, batch_size=64, patience=30,
                       use_mixup=True):
    """Train HAF-STD Net on 4D per-frame per-layer features (N, layers, T, D)."""
    model = model.to(device)
    opt = torch.optim.AdamW(model.parameters(), lr=lr, weight_decay=1e-4)
    scheduler = torch.optim.lr_scheduler.CosineAnnealingLR(opt, T_max=epochs)

    Xt = torch.tensor(X_tr_4d, dtype=torch.float32)
    yt = torch.tensor(y_tr, dtype=torch.long)
    Xe = torch.tensor(X_ev_4d, dtype=torch.float32).to(device)
    ye = torch.tensor(y_ev, dtype=torch.long).to(device)

    dl = DataLoader(TensorDataset(Xt, yt), batch_size=batch_size, shuffle=True)
    best_acc, best_state, wait = 0, None, 0

    for epoch in range(epochs):
        model.train()
        for xb, yb in dl:
            xb, yb = xb.to(device), yb.to(device)
            if use_mixup and xb.size(0) > 1:
                lam = max(np.random.beta(0.2, 0.2), 0.5)
                idx = torch.randperm(xb.size(0))
                xb_mix = lam * xb + (1 - lam) * xb[idx]
                logits = model(xb_mix)
                loss = lam * F.cross_entropy(logits, yb, label_smoothing=LABEL_SMOOTHING) + \
                       (1 - lam) * F.cross_entropy(logits, yb[idx], label_smoothing=LABEL_SMOOTHING)
            else:
                loss = F.cross_entropy(model(xb), yb, label_smoothing=LABEL_SMOOTHING)
            opt.zero_grad(); loss.backward(); opt.step()
        scheduler.step()

        model.eval()
        with torch.no_grad():
            acc = (model(Xe).argmax(1) == ye).float().mean().item()
        if acc > best_acc:
            best_acc = acc
            best_state = {k: v.cpu().clone() for k, v in model.state_dict().items()}
            wait = 0
        else:
            wait += 1
            if wait >= patience: break
    model.load_state_dict({k: v.to(device) for k, v in best_state.items()})
    return model
\end{lstlisting}

\subsection{Memory Management and Checkpointing}

Given the limited GPU memory (15.6 GB) and the sequential nature of feature extraction from two large models, careful memory management is critical. The implementation employs several strategies:

\begin{enumerate}[leftmargin=*]
    \item \textbf{Sequential Model Loading:} Only one model is loaded at a time. After V-JEPA 2 extraction is complete, the model and processor are explicitly deleted, GPU cache is cleared (\texttt{torch.cuda.empty\_cache()}), and Python garbage collection is triggered (\texttt{gc.collect()}) before loading Qwen2.5-VL.
    \item \textbf{Periodic Checkpointing:} Features and labels are saved to disk at regular intervals (every 100 clips for FineGym, every 500 videos for HMDB51). If the Kaggle session times out or the kernel restarts, extraction resumes from the last checkpoint rather than restarting.
    \item \textbf{Periodic Cache Clearing:} \texttt{torch.cuda.empty\_cache()} is called periodically (every 200--500 videos) to release fragmented GPU memory and prevent out-of-memory errors during long extraction runs.
    \item \textbf{OOM Error Recovery:} \texttt{RuntimeError} exceptions containing ``out of memory'' are caught, GPU cache is cleared, and extraction continues with the next video rather than terminating the entire pipeline.
\end{enumerate}

\subsection{Data Pipeline Challenges}

Several implementation challenges were encountered and resolved:

\begin{itemize}[leftmargin=*]
    \item \textbf{FineGym Dataset Acquisition:} A balanced dataset of 1,200 FineGym clips (300 clips per event category) was acquired using \texttt{yt-dlp} with temporal trimming based on annotation timestamps.
    \item \textbf{HMDB51 Schema Validation Error:} The Hugging Face \texttt{datasets} library encountered a \texttt{DatasetGenerationError} when loading the HMDB51 dataset through its standard API. This was resolved by bypassing the dataset library entirely and using \texttt{hf\_hub\_download} to directly download the raw zip archive, followed by manual extraction and directory-based indexing.
    \item \textbf{Qwen2.5-VL Memory Constraints:} At full float16 precision, Qwen2.5-VL (7B parameters) requires approximately 14 GB of VRAM, leaving insufficient memory for activations and intermediate computations on the 15.6 GB T4 GPU. The 4-bit NF4 quantization with double quantization reduces the model footprint to approximately 4 GB, providing adequate headroom for inference.
\end{itemize}

\subsection{Execution Timeline}

\begin{table}[H]
\centering
\caption{Execution timeline for both dataset experiments.}
\label{tab:execution_timeline}
\begin{tabular}{llr}
\toprule
\textbf{Dataset} & \textbf{Stage} & \textbf{Duration} \\
\midrule
\multirow{4}{*}{FineGym} & Video Download (\texttt{yt-dlp}) & $\sim$110 min \\
 & V-JEPA 2 (1,200 clips) & $\sim$12 min \\
 & Qwen2.5-VL (1,200 clips) & $\sim$14 min \\
 & Classification (9 classifiers) & $<$1 min \\
\midrule
\multirow{4}{*}{HMDB51} & Dataset Download (HF Hub) & $\sim$1 min \\
 & V-JEPA 2 (6,754 videos) & $\sim$32 min \\
 & Qwen2.5-VL (6,754 videos) & $\sim$56 min \\
 & Classification (9 classifiers) & $<$5 min \\
\midrule
\textbf{Total} & & \textbf{$\sim$210 min} \\
\bottomrule
\end{tabular}
\end{table}

Qwen2.5-VL takes approximately $1.75\times$ longer than V-JEPA 2 for feature extraction on HMDB51, despite processing fewer frames per video (4 vs. 8). This is attributable to the larger model size and the overhead of chat template construction and tokenization at each step. The classification stage (all 9 classifiers) adds minimal overhead, with sklearn classifiers completing in seconds and neural network classifiers requiring 1--3 minutes for training with early stopping.
